%% ==============================================================================
%% frontmatter/abstract.tex
%% Résumés bilingues (Français et Anglais) et mots-clés
%% ==============================================================================
\addchap{Résumé / Abstract}

% --- Résumé en français ---
\section*{Résumé}

La conception de documents scientifiques volumineux, tels que les mémoires de thèse et les rapports d'expertise,
se heurte fréquemment à des difficultés d'évolutivité, de cohérence graphique et de portabilité collaborative.
Ce travail propose une architecture éditoriale modulaire et découplée,
bâtie sur les fondations de la classe \LaTeX{} moderne \code{scrbook} (KOMA-Script).
En instaurant une séparation stricte entre l'orchestration du document maître (\code{main.tex}),
la classe racine (\code{mythesis.cls}),
les profils cibles écran et impression via les fichiers d'options (\code{.clo}),
ainsi que la configuration fine des styles bibliographiques (\code{.bbx}/\code{.cbx}),
cette approche garantit une clarté conceptuelle maximale.
Le système démontre qu'une structure hautement granulaire favorise le travail collaboratif sous contrôle de version,
réduit considérablement la dette technique des préambules monolithiques
et assure une conformité typographique optimale quel que soit le canal de diffusion.

\vspace{0.8em}
\noindent\textbf{Mots-clés :} \getkeywords, classe scrbook, modularité, typographie, Bib\LaTeX.

\vspace{1.5cm}

% --- Abstract en anglais ---
\begin{otherlanguage}{english}
\section*{Abstract}

The authoring of large-scale scientific manuscripts, such as doctoral dissertations and technical monographs,
often suffers from maintainability bottlenecks, visual inconsistencies, and cumbersome collaborative workflows.
This research introduces a decoupled and highly modular \LaTeX{} document architecture,
engineered on top of the modern \code{scrbook} class from the KOMA-Script suite.
By establishing a strict separation of concerns between master document orchestration (\code{main.tex}),
core document class semantics (\code{mythesis.cls}),
target-specific layout options via dedicated class options (\code{.clo}),
and fine-grained citation and bibliography tailoring (\code{.bbx}/\code{.cbx}),
this framework achieves superior transparency and maintainability.
The proposed layout demonstrates that modular decomposition mitigates merge conflicts in version-controlled repositories,
eliminates monolithic preamble pollution,
and delivers publication-grade typographical quality for both digital interactive reading and physical print production.

\vspace{0.8em}
\noindent\textbf{Keywords :} Scientific publishing, KOMA-Script, modular document architecture, \LaTeX{}, Bib\LaTeX.
\end{otherlanguage}

